\documentclass[11pt,a4paper]{article}
\usepackage{amsmath,amssymb,amsthm}
\usepackage{geometry}
\geometry{margin=1in}
\usepackage{hyperref}
\usepackage{enumitem}

\newtheorem{problem}{Problem}[section]
\theoremstyle{remark}
\newtheorem{remark}[problem]{Remark}
\newtheorem{clarification}[problem]{Clarification}
\newtheorem{motivation}[problem]{Motivation}

\newcommand{\cL}{\mathcal{L}}
\newcommand{\cH}{\mathcal{H}}
\newcommand{\cB}{\mathcal{B}}
\newcommand{\cC}{\mathcal{C}}
\newcommand{\bbC}{\mathbb{C}}
\newcommand{\bbN}{\mathbb{N}}
\newcommand{\bbR}{\mathbb{R}}
\DeclareMathOperator{\spec}{spec}
\DeclareMathOperator{\tr}{tr}
\DeclareMathOperator{\detn}{det}

\title{\bfseries Clarifications of the Problem Statements\\
\large Spectral Analysis of the Limiting Hilbert--Pólya Operator}
\author{}
\date{}

\begin{document}
\maketitle

\begin{abstract}
This document supplements the research programme \emph{Spectral Analysis of the Limiting Hilbert--Pólya Operator} by expanding on the precise content, motivation, and expected challenges of each open problem. The clarifications are intended to make the programme accessible to researchers entering the project and to forestall ambiguities that could derail the work. We explain the technical context, delineate what constitutes a solution, and point out the known pitfalls.
\end{abstract}

\section{Introduction}
The parent programme \cite{programme-H} outlines seven problems whose solution would yield a self‑adjoint operator $H$ with spectrum equal to the set of ordinates of the non‑trivial zeros of $\zeta(s)$. These problems build on the completed harmonic--categorical framework \cite{findings} and address the remaining analytic and operator‑theoretic obstacles. The present document provides extended remarks on each problem, specifying the mathematical objects involved, the type of proof expected, and the relationship to the broader Riemann Hypothesis.

Throughout we use the notation established in the original programme: $\cL_s$ is the greedy transfer operator on $H^2(\mathbb{D})$, $A(t)=\cL_{1/2+it}$, and $H_k$ are the finite‑dimensional approximations whose spectral measures converge to the zero‑counting measure on the critical line.

\section{Problem 1: Direct limit of truncated operators}

\begin{problem}[Restated]
Construct a Hilbert space $\cH_\infty$ as an inductive limit of the subspaces $\cH_k$, and define a self‑adjoint operator $H$ on $\cH_\infty$ such that the finite‑dimensional projections of $H$ reproduce the $H_k$ in the sense of spectral measures. Prove that the resolvent $(H-z)^{-1}$ is compact.
\end{problem}

\begin{clarification}
The subspaces $\cH_k$ are the domains of the truncated transfer operators $\cL_s^{(k)}$. Each $\cH_k$ is a finite‑dimensional subspace of $H^2(\mathbb{D})$, typically the span of the first $d_k$ monomials or the span of wavelet functions up to generation $k$. The inductive limit $\cH_\infty = \varinjlim \cH_k$ is an incomplete inner product space; its completion is a separable Hilbert space. 

The operator $H$ should be defined as the limit of the $H_k$ in the \emph{strong resolvent sense}: there exist isometric embeddings $V_k:\cH_k\to\cH_\infty$ such that $V_k H_k V_k^*$ converges to $H$ in the strong resolvent topology. The convergence of spectral measures already proved in \cite{findings} gives the vague convergence $\mu_k\to\mu$; the task here is to lift this to an operator convergence.

Compactness of the resolvent is essential: it guarantees that the spectrum of $H$ is purely discrete, matching the discrete set of Riemann zeros. A successful construction will produce an $H$ whose eigenvalues (with multiplicities) are exactly the limits of the eigenvalues of the $H_k$, i.e.\ the zero ordinates.
\end{clarification}

\begin{motivation}
This problem provides the most direct path from the finite approximations to a limiting Hilbert--Pólya operator. If solved, it would yield an explicit, concrete operator on a well‑understood function space, making spectral analysis straightforward.
\end{motivation}

\section{Problem 2: Generator from the unitary group}

\begin{problem}[Restated]
Recall that the problem of defining $U_t = \cL_{1/2+it}\,\cL_{1/2-it}^{-1}$ was obstructed at the zeros. However, on the orthogonal complement of the eigenspaces where $1$ is an eigenvalue, $U_t$ exists. Construct a maximal sectorially defined unitary group and use it to obtain a self‑adjoint generator $H$ via a suitable boundary‑value construction. Prove that the complementary spectrum of $H$ corresponds exactly to the Riemann zeros.
\end{problem}

\begin{clarification}
For $t$ not equal to a zero ordinate $\gamma$, the operator $A(t)=\cL_{1/2+it}$ is invertible on $\cH$ (or at least has a bounded inverse on the orthogonal complement of its kernel). Thus $U_t = A(t)A(t)^*{}^{-1} = A(t)A(-t)^{-1}$ is a well‑defined unitary on $\cH$. The family $\{U_t\}$ is defined on the open set $\bbR\setminus Z$, where $Z=\{\gamma:\zeta(1/2+i\gamma)=0\}$, and satisfies the group law wherever products are defined.

The problem asks to extend this \emph{local} unitary group to a \emph{global} object. One approach: consider the Cayley transform of $U_t$ to obtain a family of self‑adjoint operators, study their behaviour near the singular points, and glue the local generators into a single self‑adjoint operator on a larger Hilbert space that incorporates the zero data. Another approach: interpret the zeros as points where the unitary group has a ``defect'', and use boundary value theory of analytic functions to construct the generator on a space of boundary values.

A solution must (i) give a precise definition of the extended Hilbert space, (ii) define $H$ as a self‑adjoint operator on it, and (iii) prove that the point spectrum of $H$ contains $Z$ with correct multiplicities, and that the continuous spectrum is absent.
\end{clarification}

\begin{motivation}
This problem directly addresses the obstacle identified in Problem~9 of the original programme. Its solution would tie the limiting operator $H$ firmly to the transfer operator family, showing that $H$ is not an arbitrary limit but the canonical generator of the dynamical system.
\end{motivation}

\section{Problem 3: Dirac operator on the tree}

\begin{problem}[Restated]
Using the spin structure and Laplacian on the greedy tree, define a Dirac operator $D$ and a twisted version $D_s$ that encodes the transfer operator weights. Show that the spectrum of the square $D_{1/2}^2$ (or a related self‑adjoint operator) is purely discrete and equals $\{\gamma^2:\zeta(1/2+i\gamma)=0\}$.
\end{problem}

\begin{clarification}
The greedy tree $\mathcal{T}$ has a boundary $\partial\mathcal{T}=[0,\pi]$ with Lebesgue measure. The space of square‑integrable functions on the boundary is the natural Hilbert space. On the tree itself, one defines a graph Dirac operator $D:\ell^2(V)\oplus\ell^2(E)\to\ell^2(V)\oplus\ell^2(E)$ by
\[
D = \begin{pmatrix} 0 & \partial \\ \partial^* & 0 \end{pmatrix},
\]
where $\partial$ is the coboundary operator mapping vertex functions to edge functions. The square $D^2$ is block‑diagonal with the Laplacians on vertices and edges.

To incorporate the transfer operator weights, one deforms the metric on the edges using the harmonic fractions $\alpha_n$. Specifically, assign to an edge at generation $n$ the weight $w_n = \log(n+2)$. Then define a deformed Dirac operator $D_s$ whose square has eigenvalues determined by the solutions of $\cL_s\psi=\psi$. The claim is that at $s=1/2$, the spectrum of $D_{1/2}^2$ (or its positive part) consists exactly of the squares of the zero ordinates.

The problem is closely related to the spectral triple constructions of Connes and Consani \cite{ConnesConsani} and requires a detailed comparison of the tree geometry with the ad\`ele class space.
\end{clarification}

\begin{motivation}
A successful solution would give a \emph{geometric} origin for the Riemann zeros: they would arise as the spectrum of a natural differential operator on a non‑commutative space, confirming the Connes programme within the harmonic sine framework. This is conceptually the deepest of the seven problems.
\end{motivation}

\section{Problem 4: Spectral flow and holonomy}

\begin{problem}[Restated]
Study the spectral flow of the family $\cL_{1/2+it}$ along the critical line. Define the holonomy unitary on the determinant bundle and derive a self‑adjoint operator from its logarithm. Prove that the holonomy eigenvalues are exactly the zero ordinates, with multiplicities given by the winding numbers of the eigenvalues of $\cL_{1/2+it}$.
\end{problem}

\begin{clarification}
As $t$ varies, the eigenvalues $\lambda_j(t)$ of $A(t)$ move in the complex plane. The spectral flow refers to the net number of eigenvalues crossing a fixed value (here, $1$) as $t$ moves through an interval. Each zero $\gamma$ corresponds to an eigenvalue $\lambda(t)$ reaching $1$ at $t=\gamma$.

The determinant bundle over the parameter space has a natural connection induced by the family $A(t)$. The holonomy of this connection around a closed loop (or more precisely, the parallel transport between endpoints) is a unitary operator on the determinant line. Its logarithm yields a self‑adjoint operator whose eigenvalues are the arguments of the holonomy, which in turn correspond to the zero ordinates.

Concretely, one defines the function $\phi(t) = \arg\det(I-A(t))$. Its derivative $\phi'(t)$ is a spectral shift function; its points of increase are exactly the zeros. The task is to realise $\phi'(t)$ as the trace of a spectral measure of a self‑adjoint operator $H$, i.e.\ $\phi'(t) = \tr(\delta(H-t))$ in an appropriate sense.
\end{clarification}

\begin{motivation}
This approach is well‑established in the spectral theory of operator pencils and has been successful in related contexts (e.g.\ the spectral shift function of Krein). Its advantage is that it works directly on the parameter space $\bbR$, without constructing an inductive limit.
\end{motivation}

\section{Problem 5: Limit spectrum is exactly the critical zeros}

\begin{problem}[Restated]
Prove that every accumulation point of the eigenvalues of $H_k$ is either an ordinate of a critical zero or a limit of such ordinates. Show that no additional real numbers can appear in the limit spectrum.
\end{problem}

\begin{clarification}
The vague convergence $\mu_k\to\mu$ already established tells us that the limit of the spectral measures is supported on the union of the critical zero ordinates and their accumulation points. However, it does not preclude the possibility that spurious eigenvalues of $H_k$ accumulate somewhere outside the zero set in the limit. 

This problem demands a \emph{uniform} estimate: there exists a constant $C$ such that for any interval $I\subset\bbR$ not containing a zero ordinate, the number of eigenvalues of $H_k$ in $I$ is bounded independently of $k$. Equivalently, the spectral measures $\mu_k$ are \emph{tight} away from the zero set. The proof will require detailed control of the truncation error in the transfer operators and the behaviour of the entire factor $e^{P(s)}$ on the critical line.

A successful proof would show that the limiting $H$ (once constructed) has \emph{no} spectrum other than the critical zeros, without assuming the Riemann Hypothesis. The RH would then be equivalent to the statement that these zeros account for all non‑trivial zeros of $\zeta(s)$.
\end{clarification}

\begin{motivation}
This problem ensures that the Hilbert--Pólya operator does not acquire artificial spectrum from the approximation process. It is a crucial technical step that validates the entire construction.
\end{motivation}

\section{Problem 6: Connection with the explicit formula}

\begin{problem}[Restated]
Derive an explicit formula for the trace of $(H-z)^{-1}$ in terms of the logarithmic derivative of $\zeta(s)$ and the entire factor $P(s)$. Use this to prove that the spectral measure of $H$ is exactly the zero‑counting measure, without assuming RH.
\end{problem}

\begin{clarification}
The explicit formula of Guinand--Weil states, roughly, that for a test function $\phi$,
\[
\sum_{\gamma}\widehat{\phi}(\gamma) = \text{arithmetic terms} + \text{local terms at primes}.
\]
If $H$ is a self‑adjoint operator with $\spec(H)=\{\gamma\}$, then the left‑hand side is $\tr(\widehat{\phi}(H))$. The task is to identify the right‑hand side with the trace of an operator‑theoretic expression involving $\cL_s$.

More precisely, one expects an identity of the form
\[
\tr\big((H-z)^{-1}\big) = \frac{d}{dz}\log\det(I-\cL_{1/2+iz}) + \text{regular terms},
\]
which directly relates the resolvent of $H$ to the Fredholm determinant of the transfer operator. The proof will combine the determinant identity \eqref{eq:det} with the explicit formula already derived in the harmonic sine framework (Problem~2 of the original programme).

A successful proof gives an unconditional identification of the spectral measure of $H$ with the zero‑counting measure, completing the spectral interpretation independently of any unproven hypotheses.
\end{clarification}

\begin{motivation}
This problem provides the analytic link between the operator $H$ and the arithmetic data. It is the strongest possible validation that $H$ genuinely encodes the Riemann zeros and not some other set that happens to have the same distribution.
\end{motivation}

\section{Problem 7: Numerical verification and refinements}

\begin{problem}[Restated]
Implement a large‑scale numerical computation of $H_k$ for increasing $k$ and verify the convergence of the spectral measures. Investigate the statistics of eigenvalue spacings and compare with predictions from random matrix theory.
\end{problem}

\begin{clarification}
The matrices $H_k$ are explicitly computable from the truncated transfer operators. Their size $d_k$ grows with $k$. For moderate $k$ (e.g.\ $d_k\sim 1000$), standard numerical linear algebra can diagonalise $H_k$ and extract its eigenvalues. The problem asks for:
\begin{enumerate}
    \item A systematic computation of the first several hundred eigenvalues of $H_k$ for a range of $k$, to observe their convergence toward the known Riemann zeros.
    \item A statistical analysis of the eigenvalue spacings: the distribution of normalised gaps between consecutive eigenvalues should follow the GUE (Gaussian Unitary Ensemble) distribution, as predicted by the Montgomery--Odlyzko law.
    \item An investigation of the ``rate of convergence'': how fast do the eigenvalues of $H_k$ approach their limits as $k$ increases?
    \item A study of the numerical stability and the possibility of extrapolating to $k\to\infty$ using Richardson‑type acceleration.
\end{enumerate}
This problem is computational in nature but has theoretical implications: the observed GUE statistics would strongly support the conjectured self‑adjointness and the connection to random matrix theory.
\end{clarification}

\begin{motivation}
Numerical evidence is not a proof, but it provides essential guidance. If the $H_k$ eigenvalues converge quickly to the known zeros and exhibit GUE statistics, confidence in the whole framework is greatly increased. Moreover, numerical exploration may reveal patterns or unexpected properties that suggest new theorems.
\end{motivation}

\section{Interdependencies}
The problems are not ordered strictly linearly. Problems~1 and~2 are two alternative routes to construct $H$; solving either one would give a concrete operator. Problem~4 is a variant of Problem~2 with a more geometric flavour. Problem~3 is largely independent and can be pursued in parallel. Problem~5 relies on the estimates developed in Problems~2--4 but not on the full construction of $H$. Problem~6 depends on Problem~1 or~2 (one needs an operator to take the trace of). Problem~7 can begin immediately and provide data to guide the others.

\begin{thebibliography}{9}
\bibitem{programme-H} V.~Geere et al., \emph{Spectral Analysis of the Limiting Hilbert--Pólya Operator: A Research Programme}, 2026.
\bibitem{findings} ---, \emph{A Harmonic–Categorical Operator for the Riemann Zeta Function: Findings of the Research Programme}, 2026.
\bibitem{Geere2026a} V.~Geere, \emph{A Purely Geometric Reconstruction of the Sine Function}, \url{https://victorgeere.co.za/maths/sine-harmonic.html}.
\bibitem{Geere2026b} V.~Geere, \emph{On the Correlation Structure of a Greedy Harmonic Decomposition of the Sine Function}, \url{https://victorgeere.co.za/maths/greedy-harmonic-decomposition.html}.
\bibitem{Geere2026c} V.~Geere, \emph{The Quantum‑Egyptian Functional Equation}, v8, \url{https://victorgeere.co.za/maths/quantum-egyptian-functional-equation-v8.html}.
\bibitem{Mayer} D.~H.~Mayer, \emph{The thermodynamic formalism approach to Selberg's zeta function for $\mathrm{PSL}(2,\mathbb{Z})$}, Bull. Amer. Math. Soc. (N.S.) 25 (1991), 55--60.
\bibitem{Efrat} I.~Efrat, \emph{Dynamics of the continued fraction map and the spectral theory of $\mathrm{SL}(2,\mathbb{Z})$}, Invent. Math. 114 (1993), 207--218.
\bibitem{ConnesConsani} A.~Connes, C.~Consani, \emph{Schemes over $\mathbb{F}_1$ and zeta functions}, Compos. Math. 146 (2010), 1383--1415.
\end{thebibliography}

\end{document}